\section{The generalized Cayley side: information before the Borel quotient}
The trace-projection map in this program is the generalized Cayley map studied by Kostant--Michor. Their work relates its degree to the induced map on adjoint quotients and uses its differential in a harmonic-separation problem \cite{KM}. The geometric branch of the present program retains positive-degree global invariant functions, which the Borel map kills. It therefore studies additional information in the same original representation-dependent morphism.

\subsection{Quotient coordinates and the full differential}
For basic invariants $f_1,\ldots,f_r\in\C[\mathfrak t]^W$, the functions $f_i(\Theta_\rho(t))$ describe $T/W\to\mathfrak t/W$. Fundamental characters give affine coordinates on $T/W$ for a simply connected semisimple group \cite{Steinberg}. The exceptional API implements exact higher coordinates and checks invariant degrees and independence, rather than identifying a few low-degree expressions with a complete invariant system.

In logarithmic torus coordinates, the derivative matrix is
\[
 J_T(t)=M_\rho^{-1}\sum_\mu m_\mu\mu\mu^t e^\mu(t).
\]
This is only the Cartan block. Equivariance determines each root-direction eigenvalue, yielding the full left-trivialized Cayley determinant
\[
 \Psi_\rho|_T=\det J_T\prod_{\alpha>0}
 e^\alpha\left(\frac{\alpha(\Theta_\rho)}{e^\alpha-1}\right)^2.
\]
The apparent quotients extend as Laurent polynomials, because they are entries of the regular differential. The formula initially follows on a dense torus open and then extends. It is a direct deduction from the explicit torus map and equivariance. Its zero set is a \emph{source} critical locus; it is not already a computation of the target branch or nonproperness locus.

The same equivariance permits reconstruction of Borel-coordinate formulas from torus data and root commutators. Character-based recurrences extend the initial $\SL_3$ and $\Sp_4$ calculations to other representations and root data. Their value is a uniform mechanism for root coordinates and collapse tests. Singular root coefficients mark places where unipotent directions can cease to be determined by first-order data; those loci require the higher recurrence, not a generic division by a root coefficient.

\subsection{Generic degrees and boundary corrections}
For a dominant generically finite map, degree means the characteristic-zero extension degree of rational-function fields. In this setting
\[
 \deg\theta_\rho=\deg(T\to\mathfrak t)=\deg(T/W\to\mathfrak t/W).
\]
The two Weyl quotient degrees cancel; one does not divide the torus degree by $|W|$. For an adjoint character the Newton calculation often takes place on the root-lattice torus, so a central covering factor must be restored for the simply connected group.

The recorded exceptional calculations are:
\begin{center}
\begin{tabular}{lrrrr}
\toprule
Representation & Volume & Boundary & Cover & Degree\\\midrule
E6 minimal, $27$ &1080&0&1&1080\\
E6 adjoint, $78$ &2808&1890&3&2754\\
E7 minimal, $56$ &9792&2520&1&7272\\
E7 adjoint, $133$ &21312&4480&2&33664\\\bottomrule
\end{tabular}
\end{center}
The degree is $(\text{volume}-\text{boundary})\times\text{cover}$. A raw Newton volume is wrong in three of these cases. The E6 correction $1890$ also has a deformation reconstruction, supplying a different check on the intersection calculation. F4 minimal degree $96$, F4 minimal/adjoint mixtures, and defining/adjoint mixtures in B/C have separate recorded boundary formulas and exceptional coefficient ratios. These extend the method to coherent families, but do not constitute an arbitrary-representation boundary algorithm.

The unresolved structural issue is how to derive all contributing face schemes and their normal multiplicities from representation data. A facet character alone may not determine the correction. Toric intersection theory is the published foundation; the representation-specific boundary calculations are the additional work. These degrees are degrees of particular regular maps, not the minimum Cayley degree over all dominant rational equivariant maps. Historical minimum-degree questions therefore cannot be declared solved by this table.

\status{See \record{SPRINGER_GEOMETRY.md}{quotients and differential}, \record{EXCEPTIONAL_DEGREES.md}{exceptional degrees}, \record{GENERAL_CAYLEY_DEGREE.md}{boundary criteria}, \record{F4_MIXTURE_DEGREES.md}{F4 mixtures}, and \record{BC_MIXTURE_DEGREES.md}{B/C mixtures}. Exact certificate checks support the stated cases; external review and literature priority remain separate.}

\section{Regularity preservation and harmonic separation}
An element is regular if its centralizer has dimension $\rank G$. Torus regularity means regular semisimple points of $T$ map to regular points of $\mathfrak t$. Global regularity preservation also includes regular nonsemisimple group elements. These are distinct conditions, and neither is the same as generic smoothness of a specified fiber family.

The repository has a uniform classification for actual nontrivial finite-dimensional representations of simple simply connected complex groups, allowing direct sums and trivial summands. For $\SL_N$, $N\ge3$, the globally preserving representations are copies of the defining module, or copies of its dual, with optional trivial summands; mixing the two nontrivial types is excluded. For $\SL_2$, only copies of the defining module plus trivial summands preserve global regularity. No global positive cases occur in the remaining simple types under these hypotheses.

There is a larger torus-preserving symplectic family. Write $D(g)=\det(I+g)$ in the defining symplectic module. Actual characters of the form $A D^m+B$, with $A,m$ positive integers and $B$ chosen so the expression is an actual character, preserve torus regularity. In rank at least two they fail global preservation. The same torus family occurs for $\SL_2=\Sp_2$, but global preservation is more restrictive. The classification does not cover arbitrary virtual characters with negative nontrivial multiplicities, positive characteristic, or general reductive products.

The argument exploits positivity of actual weight multiplicities on real root directions, Laurent factorization of root-coordinate functions, and compatibility across the root system. Positivity is a hypothesis with content, not a property of arbitrary invariant potentials. The standard defining $\SL_N$ example is already treated by Kostant--Michor; the broader classification is a working deduction and requires focused external review.

\subsection{A sharp localization that does not require global regularity}
Let $G=\Sp(2n)$, $n\ge2$, and let $\theta_m$ be the normalized map for a determinant-power character. Put $A_G=\C[G]$, $R_G=A_G^G$, and $\mathcal H_m=\theta_m^*\operatorname{Harm}(\mathfrak g)$. The harmonic space is a vector space, not a subalgebra. The recorded theorem is
\[
 R_G[D^{-1}]\otimes_\C\mathcal H_m\xrightarrow{\sim} A_G[D^{-1}].
\]
Among nonzero invariant principal localizations by $f$, this multiplication map is an isomorphism exactly when $D$ divides $f$. Thus the determinant divisor is the minimal support to remove within that specified class of localizations. The full normalized Jacobian is
\[
 \Psi_m(g)=\left(\frac{D(g)}{4^n}\right)^{(nm-1)(2n+1)}
 \left(1-nm+\frac m2\tr(g)\right).
\]
Separation holds across the extra critical divisor in the last factor. This refines the needed localization for this fixed pulled-back harmonic space; it does not contradict the global regularity obstruction or remove poles from every inverse-differential construction. The proof uses the classical Cayley isomorphism on $D\ne0$ and known Lie-algebra harmonic separation.

\status{See \record{REGULARITY_CLASSIFICATION.md}{the complete classification}, \record{SYMPLECTIC_HARMONIC_LOCALIZATION.md}{the sharp localization}, and \record{CAYLEY_LITERATURE_AUDIT.md}{the literature audit}. These results concern the generalized Cayley map itself. They do not propose an unrelated replacement harmonic complement.}

\section{Actual adjoint fibers: the completed case theorem}
Let $G$ be simple and simply connected, $\mathfrak g=\Lie(G)$, and $B_{\rm Kill}$ its Killing form. Define
\[
 B_{\rm Kill}(\theta_{\mathrm{Ad}}(g),X)
 =\tr(\Ad(g)\operatorname{ad}X).
\]
For a proper Levi subgroup $K$ with Lie algebra $\mathfrak k$, put
\[
 Z_K^\circ=\{h\in\mathfrak z(\mathfrak k):Z_G(h)=K\}.
\]
\begin{theorem}[Generic adjoint regular-source fibers; completed repository case proof]
For every such $G,K$, there is a dense Zariski open $U_K\subset Z_K^\circ$ such that the actual scheme
\[
 \theta_{\mathrm{Ad}}^{-1}(h)\cap G_{\reg}
\]
is smooth for every $h\in U_K$. Empty fibers and smooth positive-dimensional transverse factors are allowed. After further shrinking, this is a smooth family over $U_K$.
\end{theorem}
The statement includes all finite simple types, including E8. It retains nilpotents in the equations and regular nonsemisimple source points. It does not replace a fiber by its reduced support. It does not assert that every centralizer-exact target is good, or that the fiber is nonempty, or that its dimension always equals $\dim K-\rank G$.

\subsection{The common scheme reduction}
Equivariance gives $(\Ad(g)-1)\theta(g)=0$. On the fiber over $h$, this places the centralizer equations in the fiber ideal, so the actual fiber lies in $K$. On a finite central etale cover $[K,K]\times Z(K)^\circ\to K$, use fundamental characters $x_1,\ldots,x_\ell$ of the derived group and central torus coordinates $z_1,\ldots,z_s$. Let $\Phi(x,z)$ be the \emph{full} restricted adjoint character. Define
\[
 Y=V(\partial_{x_1}\Phi,\ldots,\partial_{x_\ell}\Phi),\qquad
 \mu_j=z_j\partial_{z_j}\Phi.
\]
The transverse scheme over a target $c$ is cut out by these equations and $\mu_j=c_j$. On the regular source, the adjoint quotient is smooth, and the completed local ring of the group fiber is the completed transverse ring with $\dim K-\rank G$ formal smooth variables adjoined. The finite central cover is etale; later Weyl trace substitutions can ramify and require separate treatment.

The uniform obstruction is therefore
\[
 \overline{\mu(\operatorname{Sing}Y)}\subsetneq\mathfrak z(\mathfrak k),
\]
where $\operatorname{Sing}Y$ means the nonsmooth locus of the actual scheme. Standard generic smoothness applies only after this image is excluded. It cannot remove a persistent nonreduced component that dominates the target.

\subsection{What a certificate means in this proof}
For exceptional cases, a certificate is finite exact algebraic evidence for a precisely stated step: original polynomial generators, ideal membership witnesses or certified leading terms, dimensions or degrees, and the hypotheses linking those data to a characteristic-zero statement. It is not a numerical sample or a second name for the theorem. A monomial dimension count checks an initial ideal once that initial ideal has been justified. A finite-prime computation needs a valid characteristic-zero implication; it does not supply one by itself.

The consolidated proof handles all 255 proper E8 standard subsets through Weyl transfers and the remaining critical models. A fresh audit rechecked transfers, original critical generators, and the passage to actual group fibers. The last degree-at-infinity argument has an explicit finite-prime membership dependency on a pinned exact engine. Native checks independently verify specified identities and monomial counts, but do not independently reproduce every engine membership. That limitation should be disclosed to a reviewer rather than hidden by the phrase ``independently checked.''

Three controls remain indispensable. A special $\Sp_4$ centralizer-exact target has a singular transverse fiber, so genericity is essential. An E6 contact example has transverse excess two, so universal quasifiniteness is false. Altered characters and an actual nonadjoint $\SL_3$ representation give persistent nilpotent behavior, so Weyl symmetry and positivity alone do not prove the adjoint theorem. The case proof is complete in its recorded scope; a uniform root-theoretic replacement is still unresolved.

\status{Precise theorem and dependencies: \record{ADJOINT_GENERIC_FIBER_THEOREM.md}{the consolidation}; \record{E8_DEGREE_TRANSFER_AUDIT.md}{the fresh E8 audit}; \record{review/adjoint-generic-fibers/README.md}{the review package}. This is a computer-assisted working proof, not external mathematical certification.}

\section{Nonemptiness, excess, and representation changes}
Smoothness allowing empty fibers leaves an important classification problem: which components actually dominate a central target stratum, and how large are their fibers? The repository derives a critical-value/stabilizer criterion and a rational-linear description of the relevant central images. These are structural tools, but finding all components directly from roots remains unresolved in general.

For actual representation characters polynomial in the adjoint character, write $\chi_\sigma=f(\chi_{\mathrm{Ad}})$ and assume the trace normalization is nondegenerate. The chain rule gives
\[
 \theta_\sigma=
 \frac{f'(\chi_{\mathrm{Ad}})}{f'(\dim G)}\theta_{\mathrm{Ad}}.
\]
On a fiber over a nonzero central target the numerator scalar is a unit. On a dominant component with constant adjoint-character value $b$, the component survives exactly when $f'(b)\ne0$. Smoothness and excess are then transported by the proved component criterion. This is not a classification of arbitrary irreducible representations; actual polynomial characters are a specified family.

\subsection{Classical block classification}
For $\SL_N$ with proper Levi block sizes $b_i$, generic nonemptiness in this polynomial-adjoint family occurs exactly when every $b_i\le2$ and at least one block is a singleton. The dimension is $\sum_i b_i(b_i-1)$. In B/C/D, write the matrix Levi pattern as $\prod_i\GL_{b_i}\times H_m$, with $\sum b_i+m=n$ and residual $H_m$ orthogonal or symplectic. Absorb $\SO_2$ into a $\GL_1$ block in D.

If some $b_i>4$, or if $m>2$ in B, $m>1$ in C, or $m>3$ in D, the regular fiber is empty over centralizer-exact targets. Within those bounds, all patterns other than the all-four-block families and two small residual exceptions have a dominant component with nonconstant adjoint character, and therefore remain generically nonempty for every admissible $f$. The exceptional patterns have explicit finite constant-value sets; the criterion is that $f'$ not vanish on all those values. The two isolated patterns are B3 with $\GL_1\times\SO_5$ and D4 with $\GL_1\times\SO_6$, both governed by $f'(0)$.

Whenever nonempty, the generic fibers are smooth of dimension
\[
 \sum_i b_i(b_i-1)+2m^2\quad\text{in B/C},\qquad
 \sum_i b_i(b_i-1)+2m(m-1)\quad\text{in D}.
\]
Thus the transverse excess is zero throughout this classical family. On $\GL_4^q$ with no residual factor, the constant value is $-q/(2q-1)$ in B and zero in C/D. In particular pure adjoint tensor powers of order at least two erase the C/D strata. The exceptional E6 contact stratum behaves differently and demonstrates why this classical zero-excess conclusion cannot be promoted to every simple group.

There is also a recorded extension of adjoint generic smoothness to semisimple groups and all Levi factors, using a separate regular zero-fiber calculation and finite central descent. It is not an automatic reductive extension, since the adjoint trace form is degenerate on a central torus. The present exposition keeps the proper-Levi simple theorem as its principal statement because that is the scope of the original conjecture and its review package.

\status{See \record{ADJOINT_NONEMPTINESS_CRITICAL_VALUES.md}{the structural criterion and zero targets}, \record{CHARACTER_COMPOSITION_FIBERS.md}{polynomial-character survival}, and \record{CLASSICAL_POLYNOMIAL_ADJOINT_BLOCKS.md}{the full classical table and boundary audit}. The latter table, rather than an inferred extrapolation from the examples here, gives the complete constant-value tests.}
